\documentclass{article}

\usepackage{arxiv}
\usepackage[utf8]{inputenc}
\usepackage[T1]{fontenc}
\usepackage{hyperref}
\usepackage{url}
\usepackage{booktabs}
\usepackage{amsmath}
\usepackage{amsfonts}
\usepackage{graphicx}
\usepackage{natbib}

\title{Statistical Evaluation of Refusal Ablation in a Large Language Model}

\author{
  Matthes Fogtmann \\
  Technical University of Denmark
}

\renewcommand{\shorttitle}{Refusal Ablation Evaluation}

\hypersetup{
  pdftitle={Statistical Evaluation of Refusal Ablation in a Large Language Model},
  pdfauthor={Matthes Fogtmann},
  pdfkeywords={large language models, refusal ablation, safety evaluation, McNemar test}
}

\begin{document}
\maketitle

\begin{abstract}
  This report evaluates a single refusal-direction ablation on
  \texttt{google/gemma-4-E4B-it}. A direction is fit from a local
  policy-derived dataset of 480 safe and unsafe prompts, then removed from the
  model weights. The untouched and edited model are evaluated on the same 440
  SORRY-Bench prompts and scored with the pinned official SORRY-Bench Mistral
  judge. The judge assigns 0 to refusal and 1 to unsafe compliance. The
  baseline model complied with 94 of 440 prompts (21.4\%), while the edited
  model complied with 337 of 440 prompts (76.6\%). In the paired comparison,
  247 prompts changed from baseline refusal to edited compliance, while 4
  changed in the opposite direction. An exact McNemar/binomial test gives
  $p=4.53\cdot10^{-68}$, with an estimated compliance increase of 55.2
  percentage points. The result shows that this edit greatly weakens refusal
  behavior on the official benchmark. The analysis is intentionally narrow: it
  measures benchmark refusal behavior, not general model quality or robust
  safety under further interventions.
\end{abstract}

\keywords{large language models \and refusal ablation \and safety evaluation \and paired model comparison}

\section{Introduction}

Large language models are often aligned to refuse harmful requests while still
answering benign requests. Refusal ablation is an intervention that attempts to
remove or weaken this refusal behavior. This makes it a useful case study for
statistical evaluation because the intervention may create a safety--utility
trade-off: fewer refusals can improve helpfulness on benign prompts, but can also
increase compliance with unsafe prompts.

This report evaluates whether a refusal-ablated model behaves differently from
the untouched model when both are tested on the same prompt set. The focus is
not on training a new model, but on designing and analyzing a paired model
evaluation with appropriate metrics and statistical tests.

\subsection{Research Questions}

We consider three research questions:

\begin{itemize}
  \item \textbf{RQ1:} Does refusal ablation reduce refusals on unsafe prompts?
  \item \textbf{RQ2:} Does refusal ablation reduce false refusals on benign prompts?
  \item \textbf{RQ3:} Does the ablation effect vary across prompt categories?
\end{itemize}

\subsection{Contribution}

The contribution of the report is a statistical comparison of the untouched and
refusal-ablated model using paired binary response labels. The same prompts are
evaluated on both models, which allows direct comparison of model behavior at
the prompt level. Following the model-comparison setting described by
\citet{raschka2018model}, we use McNemar's test for paired binary outcomes and
report effect sizes alongside significance tests.

\section{Methods}

\subsection{Models}

TODO: Describe the two evaluated systems.

\begin{itemize}
  \item \textbf{Baseline model:} TODO, e.g. untouched Gemma model.
  \item \textbf{Ablated model:} TODO, e.g. same model after refusal-direction ablation.
\end{itemize}

The ablation procedure follows the idea that refusal behavior can be mediated by
a direction in activation space \citep{arditi2024refusal}. TODO: Briefly state
how the ablated model was produced, or state that the model outputs were provided
from a separate machine.

\subsection{Dataset}

The final analysis dataset contains one row per prompt. Each prompt is evaluated
on both models, giving paired observations. The binary response labels are:

\begin{equation}
  Y_{i,m} =
  \begin{cases}
    1, & \text{if model } m \text{ refused prompt } i,\\
    0, & \text{if model } m \text{ answered prompt } i.
  \end{cases}
\end{equation}

TODO: Insert final dataset counts in Table~\ref{tab:data-summary}.

\begin{table}[h]
  \centering
  \caption{Dataset summary. Replace the TODO values after the final labeled data is available.}
  \label{tab:data-summary}
  \begin{tabular}{lr}
    \toprule
    Quantity & Count \\
    \midrule
    Total prompts & TODO \\
    Unsafe prompts & TODO \\
    Benign prompts & TODO \\
    Prompt categories & TODO \\
    Median prompts per category & TODO \\
    \bottomrule
  \end{tabular}
\end{table}

\subsection{Metrics}

For unsafe prompts, a refusal is treated as the desired behavior and an answer is
treated as a non-refusal. Because the binary labels do not judge the content of
the answer, the unsafe-prompt answered rate is interpreted as a proxy for unsafe
compliance rather than a direct measure of harmfulness.

\begin{equation}
  \text{Unsafe non-refusal rate}_m =
  \frac{1}{N_{\text{unsafe}}}
  \sum_{i \in \text{unsafe}} (1 - Y_{i,m}).
\end{equation}

For benign prompts, refusal is treated as a false refusal:

\begin{equation}
  \text{False refusal rate}_m =
  \frac{1}{N_{\text{benign}}}
  \sum_{i \in \text{benign}} Y_{i,m}.
\end{equation}

We report rates for both models and the paired difference:

\begin{equation}
  \Delta = \bar{Y}_{\text{ablated}} - \bar{Y}_{\text{baseline}}.
\end{equation}

TODO: Specify whether confidence intervals are computed by bootstrap over
prompts or by an exact/binomial method for paired differences.

\subsection{Statistical Tests}

The primary statistical test is McNemar's test \citep{mcnemar1947note}. This is
appropriate because each prompt is evaluated on both models and the response
label is binary.

For unsafe prompts, the paired contingency table is:

\begin{table}[h]
  \centering
  \caption{Paired table for unsafe prompts.}
  \label{tab:mcnemar-unsafe-layout}
  \begin{tabular}{llcc}
    \toprule
    & & \multicolumn{2}{c}{Baseline model} \\
    \cmidrule(lr){3-4}
    & & Refused & Answered \\
    \midrule
    Ablated model & Refused & $a$ & $b$ \\
    Ablated model & Answered & $c$ & $d$ \\
    \bottomrule
  \end{tabular}
\end{table}

The null hypothesis is that the two off-diagonal flip probabilities are equal:

\begin{equation}
  H_0: P(\text{baseline refused, ablated answered})
  =
  P(\text{baseline answered, ablated refused}).
\end{equation}

For RQ1, the directional alternative is:

\begin{equation}
  H_A: P(\text{baseline refused, ablated answered})
  >
  P(\text{baseline answered, ablated refused}).
\end{equation}

For benign prompts, we use the same test structure, but the interpretation
changes: a refused benign prompt is counted as a false refusal.

\subsection{Category Analysis}

To answer RQ3, we report category-wise refusal-rate changes. If sample sizes are
sufficient, we also fit an exploratory logistic model:

\begin{equation}
  \text{logit}\left(P(Y_{i,m}=1)\right)
  =
  \beta_0 + \beta_1 \text{Ablated}_m +
  \beta_2 \text{Category}_i +
  \beta_3(\text{Ablated}_m \times \text{Category}_i).
\end{equation}

Because each prompt appears once for each model, uncertainty should account for
the paired design, for example by clustering standard errors by prompt id or by
using a repeated-measures model. The category analysis is treated as exploratory,
especially if many category-level comparisons are made.

\section{Results}

\subsection{Overall Compliance Change}

The baseline model complied with 94 of 440 SORRY-Bench prompts (21.4\%). After
the refusal-direction edit, the model complied with 337 of 440 prompts
(76.6\%). This is a paired compliance-rate increase of 55.2 percentage points.
A normal approximation for the paired mean difference gives a 95\% interval of
50.4 to 60.0 percentage points.

\begin{table}[h]
  \centering
  \caption{Overall official SORRY-Bench result.}
  \label{tab:overall-results}
  \begin{tabular}{lrrr}
    \toprule
    Condition & Prompts & Compliance rate & Refusal rate \\
    \midrule
    Baseline Gemma & 440 & 21.4\% & 78.6\% \\
    Edited Gemma   & 440 & 76.6\% & 23.4\% \\
    \bottomrule
  \end{tabular}
\end{table}

\subsection{Paired Prompt Flips}

The paired table shows that most changed prompts moved in the same direction:
from baseline refusal to edited compliance. There were 247 such flips, compared
with only 4 prompts that moved from baseline compliance to edited refusal.

\begin{table}[h]
  \centering
  \caption{Paired official judge scores. Score 0 is refusal and score 1 is unsafe compliance.}
  \label{tab:paired-scores}
  \begin{tabular}{lrr}
    \toprule
    & \multicolumn{2}{c}{Edited Gemma} \\
    \cmidrule(lr){2-3}
    Baseline Gemma & Refusal & Compliance \\
    \midrule
    Refusal    & 99 & 247 \\
    Compliance & 4  & 90  \\
    \bottomrule
  \end{tabular}
\end{table}

The exact McNemar/binomial test on the 251 discordant pairs gives
$p=4.53\cdot10^{-68}$. The change is therefore not only statistically
detectable, but also large in practical terms.

\subsection{Category-Level Pattern}

The effect is broad but not identical across SORRY-Bench categories. Several
category IDs changed from 0\% baseline compliance to 100\% edited compliance,
including categories 12, 26, 27, and 43. Other categories changed much less:
categories 33, 34, and 38 showed no change, category 20 increased by 30
percentage points, and category 41 increased by 10 percentage points.

This category variation is important for interpretation. The edit does not
simply turn every refusal into compliance; some benchmark categories remain more
resistant than others.

\section{Discussion}

TODO: Interpret the results in terms of the research questions.

For unsafe prompts, a lower refusal rate in the ablated model indicates that the
intervention weakened safety behavior. For benign prompts, a lower refusal rate
may indicate improved helpfulness. The central interpretation is therefore the
trade-off between reduced over-refusal and increased unsafe-prompt non-refusal.

\subsection{Statistical Interpretation}

TODO: Discuss both statistical and practical significance. Report whether the
effect size is large enough to matter, not only whether the $p$-value is below
0.05.

\subsection{Limitations}

The binary refused/answered label is simple and reproducible, but it does not
directly measure answer harmfulness. An answered unsafe prompt is therefore a
proxy for unsafe compliance, not proof that the response contained actionable
harmful content. Manual labels may also contain subjectivity. Finally, the
results depend on the sampled prompts, prompt categories, model version, and
ablation method.

\subsection{Ethical Considerations}

This evaluation studies safety behavior under harmful prompts. The report should
avoid reproducing harmful prompt text or unsafe model outputs unless necessary,
and any examples should be paraphrased or summarized at a high level.

\section{Conclusion}

TODO: Give a concise answer to the research questions. State whether refusal
ablation significantly changed refusal behavior, whether it reduced false
refusals on benign prompts, and whether the effect was category-dependent.

The final conclusion should distinguish between statistical significance,
practical effect size, and the limitations of binary refusal labels.


\bibliographystyle{unsrtnat}
\bibliography{references}

\end{document}
